\chapter{Integer datapath for hardware}
\label{ch:hardware}

This chapter turns the trained spiking block into arithmetic that a hardware accelerator can implement with
integers only, and describes how energy is estimated. The goal is to accelerate the six spiking blocks; the
embedding, input neurons and output head stay on a host processor.

\section{Where the operations are}

Instrumenting the trained model on validation text gave the operation counts in Table~\ref{tab:ops}. Two facts
shaped the design.

\begin{table}[t]
\centering
\caption[Operations per token inside the spiking blocks]{Operations per generated token inside the six blocks of
the final model ($T=8$, 256-token context). ``Events'' are non-zero inputs, each of which adds one weight column.}
\label{tab:ops}
\small
\begin{tabular}{llr}
\toprule
Operation & Inputs & Per token \\
\midrule
Linear layers reading the residual stream (Q/K/V, first MLP layer) & integers 0--13, 58\% non-zero & 11.9M events \\
Linear layers reading spikes (projection, second MLP layer) & binary, 19\% non-zero & 3.3M events \\
Normalisation & 122{,}880 elements & 0.5M operations \\
LIF neurons & 122{,}880 neurons & 0.6M operations \\
Attention (AND + popcount, square, accumulate) & 25\% of pairs overlap & 0.13M operations \\
\midrule
Weight columns fetched (once per token for all 8 steps) & & 2.78M weights \\
\bottomrule
\end{tabular}
\end{table}

First, \textbf{the linear layers are 98\% of the arithmetic}, and most of them read the residual stream, whose
values are not spikes but small integers. A synapse that receives the value 5 must add its weight five times or
multiply. The claim that spiking networks need only additions does not hold for these layers.

Second, \textbf{memory is the main cost}. The block has 4.7 million weights. If each layer processes all 8 time
steps of a token together, a weight column can be fetched once and reused for every step in which its input is
non-zero; a column whose input never fires is not fetched at all.

\section{Shift-coded residual inputs}

Each residual value $r \in \{0, \dots, 13\}$ is replaced by its leading bit, rounded on the first bit that is cut
and saturated at 8:
\begin{center}
\begin{tabular}{l*{14}{c}}
\toprule
$r$ & 0 & 1 & 2 & 3 & 4 & 5 & 6 & 7 & 8 & 9 & 10 & 11 & 12 & 13 \\
coded & 0 & 1 & 2 & 4 & 4 & 4 & 8 & 8 & 8 & 8 & 8 & 8 & 8 & 8 \\
\bottomrule
\end{tabular}
\end{center}
The product of a weight with a power of two is a shift, so the multiplier becomes a 4-way multiplexer that selects
$w$, $2w$, $4w$ or $8w$. Cutting instead of rounding, or keeping two bits, was also tested
(Section~\ref{sec:quant}); one rounded bit was the cheapest option that cost almost nothing. Values 12 and 13 would
round to 16 but occur for only 0.001\% of inputs, so saturating at 8 is free.

\section{Integer weights and normalisation}
\label{sec:intnorm}

\paragraph{Weights.} All weights of a linear layer share one int8 scale. Because the normalisation that follows
divides by the vector's own spread, a common scale cancels: the normalisation can work directly on the integer
accumulators, and the scale never has to be applied. One scale per output channel would be slightly more accurate
but would need a multiplication per element to bring the channels to a common scale.

\paragraph{Normalisation.} For a vector $x$ of $N$ accumulator values ($N$ = 256, 1{,}024 or 32 for a head):
\begin{enumerate}
  \item mean $= \mathrm{round}(\sum x / N)$: an adder tree and a shift, since $N$ is a power of two;
  \item centred values $c = x - \text{mean}$ and $\mathrm{MAD} = \mathrm{round}(\sum |c| / N)$, held in 12 bits;
  \item the reciprocal $1/\mathrm{MAD}$ from a leading-one detector and a table of 16 entries, indexed by the four
        bits after the leading one;
  \item $n = c \times \text{table entry}$, shifted by the leading-one position and rounded to Q4.7.
\end{enumerate}
The multiplication in step 4 (19 by 6 bits) is the only multiplier left in the block. Replacing it with a single
shift, or the table with four entries, was tested and was not accurate enough (Section~\ref{sec:quant}).

\section{Folding the normalisation into the neuron}

The normalisation ends with $\gamma n + \beta$, which feeds the LIF neuron of Equation~\ref{eq:lif}. Because the
neuron is linear below its threshold, both parameters can be moved into it. Dividing the membrane potential by
$\gamma$ gives an equivalent neuron
\begin{equation}
  u \leftarrow \big(u + n + \beta/\gamma\big) \gg 1, \qquad \text{spike if } \gamma\,(u - 1/\gamma) \ge 0,
  \qquad u \leftarrow u - 1/\gamma \text{ after a spike},
\end{equation}
with three constants per neuron: $1/\gamma$, $\beta/\gamma$ and the sign of $\gamma$, which decides the direction of
the comparison. This is exact and removes the per-element multiplication by $\gamma$. All values are held in Q4.7.
For 40 of the 15{,}360 neurons $|\gamma| < 0.05$, so $1/\gamma$ would not fit in 12 bits; these neurons barely depend on
their input, fire at most once or twice per token, and are set to never fire, which does not change the accuracy.

\section{Rounding, saturation and the attention}

All rounding is half up (add half, then cut), which needs one adder. All registers saturate rather than wrap:
accumulators at 18 bits (computed in 21 bits so that the summation order does not matter), attention sums at 16
bits, the MAD at 12 bits, and neuron states at 12 bits. The residual stream never exceeds 13 and needs 4 bits. In the
attention, counts, squares and sums are exact integers; the attention's fixed scale of $1/8$ is absorbed into the head
normalisation's $\varepsilon$.

\section{Quantization-aware training and the reference model}

The trained model is fine-tuned for 5{,}000 steps with this entire datapath in the forward pass and straight-through
gradients. The neurons' parameters stay $\gamma$ and $\beta$ and are folded on the fly, so the checkpoints keep their
format.

To hand the design to hardware, a reference model was written in pure integer arithmetic (NumPy int64). It reproduces
the PyTorch simulation of the datapath exactly: 100\% of block outputs agreed over 8 test windows of 256 tokens, six
blocks and eight time steps, and the test scores were identical. Writing it also exposed two bugs in the simulation
--- on the GPU, \texttt{log2(8.0)} returned 2.999\dots{} and powers of two computed with \texttt{pow} were inexact,
which flipped some roundings --- both fixed. The reference model produces test vectors for every stage of the block,
and a specification documents every format, rounding rule and table (summarised in Appendix~\ref{app:spec}).

\section{Energy model}
\label{sec:energymodel}

Energy is estimated per generated token as the sum, over operation types, of the number of operations times the
energy of one operation. The counts are measured on the models; the energies follow the 45\,nm figures of
\citet{horowitz2014energy}, with interpolations where needed (Appendix~\ref{app:energy}). Weights and caches are
assumed to sit in on-chip memory at 4\,pJ per byte, between the cost of a 32\,KB and a 1\,MB memory; since this
assumption matters, results are also given for 1.25, 12.5 and 160\,pJ per byte. Leakage, control logic and data
movement between units are not included, so the numbers are best read as ratios between models.
